\itemhead{3.1 例题 · 例题3（教材p.67）}{XJ-B1-C03-EX-P067-Y0440}
\label{XJ-B1-C03-EX-P067-Y0440}
{\footnotesize 湘教版教材p.67 / PDF第74页 / 3.1 例题 / 例题3}
\field{原题}{已知定义域为 $\mathbb R$ 的函数 $f(x)=x+1$ 和 $g(x)=x^2$，计算下列各式：
\subq{1}$f(2)+g(3)$；
\subq{2}$f(a^2)-g(a)$；
\subq{3}$f(f(f(0)))$。}
\field{参考答案}{（1）12；（2）1；（3）3。复合求值从最内层开始代入。}
{\small \field{关联知识}{F03 函数求值、解析式与函数建模（教材帮p.80、87、88、89、90、131 / PDF89、96、97、98、99、140）}
\field{方法与易错点}{标清每次输入和输出；按已知函数类型或代换关系建立等价方程。 易错：将 $f(x)$ 当成 $f	imes x$；复合顺序反；换元后沿用错误定义域。}
\field{答案依据}{例题原解位置：湘教版教材p.67及题库原图；独立编写可编辑参考解答。答案册此部分主要收练习、习题，不伪造例题答案册位置。}}
